% Figure: two-panel regression diagnostic for the worked tensile
% fit. Left: data with fitted line yhat = 2.330 + 0.871 x.
% Right: residuals against fitted values, with a zero reference.
% Coordinates are exactly the worked-example numbers in the prose.
\begin{figure}[htbp]
\centering
\begin{tikzpicture}
\begin{axis}[
    name=fit,
    width=7.5cm, height=6cm,
    xmin=0, xmax=9, ymin=2, ymax=10,
    axis lines=left,
    xlabel={Carbon content $x$ (0.01\,\%)},
    ylabel={Yield strength $y$ (working units)},
    title={\small Fit: $\hat{y} = 2.293 + 0.907\,x$},
    every axis plot/.append style={mark=none},
]
\addplot[only marks, engblue, mark=*, mark size=1.4pt] coordinates {
  (1,3.1) (2,4.0) (3,5.2) (4,5.8) (5,7.1) (6,7.6) (7,9.0) (8,9.2)
};
\addplot[engred, very thick, domain=0.5:8.5] {2.293 + 0.907*x};
\end{axis}

\begin{axis}[
    name=resid,
    at={($(fit.east)+(1.6cm,0)$)}, anchor=west,
    width=7.5cm, height=6cm,
    xmin=2.5, xmax=10, ymin=-0.7, ymax=0.7,
    axis lines=left,
    xlabel={Fitted value $\hat{y}$},
    ylabel={Residual $\hat{\varepsilon}$},
    title={\small Residuals against fitted},
    every axis plot/.append style={mark=none},
]
\addplot[engblack, dashed] coordinates {(2.5,0) (10,0)};
\addplot[only marks, engblue, mark=*, mark size=1.4pt] coordinates {
  (3.20,-0.10) (4.11,-0.11) (5.01,0.19) (5.92,-0.12)
  (6.83,0.27) (7.74,-0.14) (8.64,0.36) (9.55,-0.35)
};
\end{axis}
\end{tikzpicture}
\caption[Worked regression fit with residual diagnostic]{The
worked tensile-strength fit. Left: the eight specimens with the
ordinary-least-squares line. Right: residuals against fitted
values, the first plot an engineer should read. A horizontal band
of roughly constant width with no curvature is consistent with the
homoscedastic Gaussian model; the worked residuals scatter around
zero without a systematic pattern, so the linear model is
defensible on these data. A funnel or a curve in this panel would
override a high $R^{2}$.}
\label{fig:vol02:ch14:regression-residuals}
\end{figure}
